% Shared preamble for the course decision diagrams.
% Each diagram file starts with % Shared preamble for the course decision diagrams.
% Each diagram file starts with % Shared preamble for the course decision diagrams.
% Each diagram file starts with \input{diagram_style} and contains one tikzpicture.
% Compile with pdflatex; build_diagrams.py then crops each page to its content.
\documentclass[11pt]{article}
\usepackage[a3paper,landscape,margin=10mm]{geometry}
\usepackage[T1]{fontenc}
\usepackage{tgheros}                 % Neudata body font (as in neudata.cls)
\renewcommand{\sfdefault}{qhv}
\renewcommand{\familydefault}{\sfdefault}
\usepackage[scaled=0.95,varqu]{zi4}
\usepackage{tikz}
\usetikzlibrary{arrows.meta,positioning,fit,calc}
\pagestyle{empty}
\setlength{\parindent}{0pt}
\hyphenpenalty=10000
\exhyphenpenalty=10000

% Neudata palette (neu-data/.github BRAND.md); names kept so the diagram sources need no edits
\definecolor{cblue}{HTML}{0B376C}    % Neudata Blue: question boxes, links between diagrams
\definecolor{corange}{HTML}{B5652B}  % muted contrasting colour: warnings
\definecolor{cgreen}{HTML}{055F56}   % Neudata Teal: method boxes
\definecolor{clight}{HTML}{EAF2F5}   % Neudata Mist: question box fill
\definecolor{ctext}{HTML}{1F2D33}    % Neudata Ink: text
\definecolor{cgrey}{HTML}{8A99A3}    % Neudata Grey: further-learning routes
\definecolor{cpale}{HTML}{F4F8FA}
\definecolor{cnavy}{HTML}{04242F}    % Deep Navy: titles

\tikzset{
  every node/.style={text=ctext, font=\small},
  title/.style={font=\Large\bfseries, text=cnavy, anchor=west},
  subtitle/.style={font=\normalsize, text=ctext!72, anchor=west},
  q/.style={draw=cblue, fill=clight, rounded corners=3pt, line width=0.8pt,
            align=center, inner sep=6pt, text width=#1},
  q/.default=5cm,
  m/.style={draw=cgreen, fill=white, rounded corners=3pt, line width=1pt,
            align=left, inner sep=6pt, text width=#1},
  m/.default=5cm,
  opt/.style={draw=cblue, fill=white, rounded corners=3pt, line width=0.6pt,
              align=center, inner sep=5pt, text width=#1},
  opt/.default=3.2cm,
  further/.style={draw=cgrey, dashed, fill=white, rounded corners=3pt, line width=0.8pt,
                  align=left, inner sep=6pt, text=ctext, text width=#1},
  further/.default=5cm,
  warn/.style={draw=corange, fill=white, rounded corners=3pt, line width=1pt,
               align=left, inner sep=6pt, text width=#1},
  warn/.default=6cm,
  note/.style={align=left, text width=#1, font=\footnotesize, text=ctext},
  note/.default=5cm,
  band/.style={font=\normalsize\bfseries, text=cgreen, anchor=west, align=left},
  flow/.style={-{Stealth[length=2.6mm]}, line width=0.8pt, draw=ctext!75},
  flowg/.style={-{Stealth[length=2.6mm]}, line width=0.8pt, draw=cgrey, dashed},
  lab/.style={font=\footnotesize, fill=white, inner sep=2pt, text=ctext},
  code/.style={font=\ttfamily\footnotesize},
}
\newcommand{\rc}[1]{{\ttfamily\footnotesize #1}}
\newcommand{\dref}[1]{{\color{cblue}#1}}
 and contains one tikzpicture.
% Compile with pdflatex; build_diagrams.py then crops each page to its content.
\documentclass[11pt]{article}
\usepackage[a3paper,landscape,margin=10mm]{geometry}
\usepackage[T1]{fontenc}
\usepackage{tgheros}                 % Neudata body font (as in neudata.cls)
\renewcommand{\sfdefault}{qhv}
\renewcommand{\familydefault}{\sfdefault}
\usepackage[scaled=0.95,varqu]{zi4}
\usepackage{tikz}
\usetikzlibrary{arrows.meta,positioning,fit,calc}
\pagestyle{empty}
\setlength{\parindent}{0pt}
\hyphenpenalty=10000
\exhyphenpenalty=10000

% Neudata palette (neu-data/.github BRAND.md); names kept so the diagram sources need no edits
\definecolor{cblue}{HTML}{0B376C}    % Neudata Blue: question boxes, links between diagrams
\definecolor{corange}{HTML}{B5652B}  % muted contrasting colour: warnings
\definecolor{cgreen}{HTML}{055F56}   % Neudata Teal: method boxes
\definecolor{clight}{HTML}{EAF2F5}   % Neudata Mist: question box fill
\definecolor{ctext}{HTML}{1F2D33}    % Neudata Ink: text
\definecolor{cgrey}{HTML}{8A99A3}    % Neudata Grey: further-learning routes
\definecolor{cpale}{HTML}{F4F8FA}
\definecolor{cnavy}{HTML}{04242F}    % Deep Navy: titles

\tikzset{
  every node/.style={text=ctext, font=\small},
  title/.style={font=\Large\bfseries, text=cnavy, anchor=west},
  subtitle/.style={font=\normalsize, text=ctext!72, anchor=west},
  q/.style={draw=cblue, fill=clight, rounded corners=3pt, line width=0.8pt,
            align=center, inner sep=6pt, text width=#1},
  q/.default=5cm,
  m/.style={draw=cgreen, fill=white, rounded corners=3pt, line width=1pt,
            align=left, inner sep=6pt, text width=#1},
  m/.default=5cm,
  opt/.style={draw=cblue, fill=white, rounded corners=3pt, line width=0.6pt,
              align=center, inner sep=5pt, text width=#1},
  opt/.default=3.2cm,
  further/.style={draw=cgrey, dashed, fill=white, rounded corners=3pt, line width=0.8pt,
                  align=left, inner sep=6pt, text=ctext, text width=#1},
  further/.default=5cm,
  warn/.style={draw=corange, fill=white, rounded corners=3pt, line width=1pt,
               align=left, inner sep=6pt, text width=#1},
  warn/.default=6cm,
  note/.style={align=left, text width=#1, font=\footnotesize, text=ctext},
  note/.default=5cm,
  band/.style={font=\normalsize\bfseries, text=cgreen, anchor=west, align=left},
  flow/.style={-{Stealth[length=2.6mm]}, line width=0.8pt, draw=ctext!75},
  flowg/.style={-{Stealth[length=2.6mm]}, line width=0.8pt, draw=cgrey, dashed},
  lab/.style={font=\footnotesize, fill=white, inner sep=2pt, text=ctext},
  code/.style={font=\ttfamily\footnotesize},
}
\newcommand{\rc}[1]{{\ttfamily\footnotesize #1}}
\newcommand{\dref}[1]{{\color{cblue}#1}}
 and contains one tikzpicture.
% Compile with pdflatex; build_diagrams.py then crops each page to its content.
\documentclass[11pt]{article}
\usepackage[a3paper,landscape,margin=10mm]{geometry}
\usepackage[T1]{fontenc}
\usepackage{tgheros}                 % Neudata body font (as in neudata.cls)
\renewcommand{\sfdefault}{qhv}
\renewcommand{\familydefault}{\sfdefault}
\usepackage[scaled=0.95,varqu]{zi4}
\usepackage{tikz}
\usetikzlibrary{arrows.meta,positioning,fit,calc}
\pagestyle{empty}
\setlength{\parindent}{0pt}
\hyphenpenalty=10000
\exhyphenpenalty=10000

% Neudata palette (neu-data/.github BRAND.md); names kept so the diagram sources need no edits
\definecolor{cblue}{HTML}{0B376C}    % Neudata Blue: question boxes, links between diagrams
\definecolor{corange}{HTML}{B5652B}  % muted contrasting colour: warnings
\definecolor{cgreen}{HTML}{055F56}   % Neudata Teal: method boxes
\definecolor{clight}{HTML}{EAF2F5}   % Neudata Mist: question box fill
\definecolor{ctext}{HTML}{1F2D33}    % Neudata Ink: text
\definecolor{cgrey}{HTML}{8A99A3}    % Neudata Grey: further-learning routes
\definecolor{cpale}{HTML}{F4F8FA}
\definecolor{cnavy}{HTML}{04242F}    % Deep Navy: titles

\tikzset{
  every node/.style={text=ctext, font=\small},
  title/.style={font=\Large\bfseries, text=cnavy, anchor=west},
  subtitle/.style={font=\normalsize, text=ctext!72, anchor=west},
  q/.style={draw=cblue, fill=clight, rounded corners=3pt, line width=0.8pt,
            align=center, inner sep=6pt, text width=#1},
  q/.default=5cm,
  m/.style={draw=cgreen, fill=white, rounded corners=3pt, line width=1pt,
            align=left, inner sep=6pt, text width=#1},
  m/.default=5cm,
  opt/.style={draw=cblue, fill=white, rounded corners=3pt, line width=0.6pt,
              align=center, inner sep=5pt, text width=#1},
  opt/.default=3.2cm,
  further/.style={draw=cgrey, dashed, fill=white, rounded corners=3pt, line width=0.8pt,
                  align=left, inner sep=6pt, text=ctext, text width=#1},
  further/.default=5cm,
  warn/.style={draw=corange, fill=white, rounded corners=3pt, line width=1pt,
               align=left, inner sep=6pt, text width=#1},
  warn/.default=6cm,
  note/.style={align=left, text width=#1, font=\footnotesize, text=ctext},
  note/.default=5cm,
  band/.style={font=\normalsize\bfseries, text=cgreen, anchor=west, align=left},
  flow/.style={-{Stealth[length=2.6mm]}, line width=0.8pt, draw=ctext!75},
  flowg/.style={-{Stealth[length=2.6mm]}, line width=0.8pt, draw=cgrey, dashed},
  lab/.style={font=\footnotesize, fill=white, inner sep=2pt, text=ctext},
  code/.style={font=\ttfamily\footnotesize},
}
\newcommand{\rc}[1]{{\ttfamily\footnotesize #1}}
\newcommand{\dref}[1]{{\color{cblue}#1}}

\begin{document}
\begin{tikzpicture}[x=1cm,y=1cm]
\node[title] at (0,2.2) {D3. Categorical outcomes: tables, tests and effect measures};
\node[subtitle] at (0,1.5) {Binary or nominal outcome compared between groups or exposures. Always state the denominator.};

\node[q=11cm, anchor=north] (root) at (9.9,0.6) {\textbf{Categorical outcome by group or exposure}\\\footnotesize e.g.\ high BP (yes/no) by smoking; adverse event by programme};
\node[m=11cm, below=5mm of root] (table) {\textbf{Start with a contingency table: counts and row percentages}\\
Row percentage = number with the outcome / row total (e.g.\ smokers with high BP / all smokers)\\
\rc{tab <- table(hearts\$smoker, hearts\$high\_bp\_0)}\\\rc{prop.table(tab, margin = 1)}};
\draw[flow] (root) -- (table);
\node[q=7cm, below=6mm of table] (design) {\textbf{Are the observations independent or paired?}};
\draw[flow] (table) -- (design);

% independent
\node[q=7.2cm, anchor=north] (exp) at ($(design.south)+(-4.1,-1.0)$) {\textbf{Independent groups}\\\footnotesize Are the expected counts large enough (about 5 or more in each cell)?};
\draw[flow] (design.south) -- ++(0,-0.45) -| node[lab, pos=0.25] {independent} (exp.north);
\node[m=5.3cm, anchor=north] (chi) at ($(exp.south)+(-2.9,-1.0)$) {\textbf{Chi-square test}\\
\rc{chisq.test(tab)}\\
\rc{chisq.test(tab)\$expected}\\
2 x 2 tables: Yates' correction by default. Works for larger r x c tables too.};
\node[m=5.3cm, anchor=north] (fis) at ($(exp.south)+(2.9,-1.0)$) {\textbf{Fisher's exact test}\\
\rc{fisher.test(tab\_ae)}\\
No large-sample approximation; usual choice when expected counts are small (rare outcomes, small groups).};
\draw[flow] (exp.south) -- ++(0,-0.45) -| node[lab, pos=0.25] {yes} (chi.north);
\draw[flow] (exp.south) -- ++(0,-0.45) -| node[lab, pos=0.25] {no, small} (fis.north);

% paired
\node[m=5.4cm, anchor=north] (mcn) at ($(design.south)+(4.9,-1.0)$) {\textbf{McNemar's test}\\
Same people, yes/no twice (high BP at baseline and 12 months), or matched pairs. Only pairs that changed carry information.\\
\rc{mcnemar.test(tab\_paired)}};
\draw[flow] (design.south) -- ++(0,-0.45) -| node[lab, pos=0.3] {paired} (mcn.north);

\node[warn=11.6cm, anchor=north west] (sig) at ($(chi.south west)+(0,-0.6)$) {\textbf{A significant test says the proportions differ somewhere.} In tables larger than 2 x 2, look at the row percentages to see where. Report the effect measure (right) with its 95\% CI, not only the p-value.};
\node[further=5.4cm, below=4mm of mcn] (ord) {Ordinal outcome (e.g.\ satisfaction 1 to 5): the chi-square test ignores the order. Consider a rank-based test (\dref{D2}) or ordinal regression (further learning).};

% effect measures by design
\node[q=8.6cm, anchor=north west] (meas) at (18.6,0.6) {\textbf{Which effect measure? It depends on the study design}};
\node[m=8.6cm, below=5mm of meas] (coh) {\textbf{Cohort study or randomised trial}\\
People followed over time, so proportions are \textbf{risks}.\\
Report risk difference (RD) and risk ratio (RR); OR also possible.\\
Trial example, adverse events: 14/240 (5.8\%) with the programme vs 3/240 (1.2\%) with standard advice. RR 4.67 (95\% CI 1.36 to 16.03).};
\node[m=8.6cm, below=4mm of coh] (cs) {\textbf{Cross-sectional study} (one time point)\\
Proportions are \textbf{prevalences}. Report prevalence difference and prevalence ratio (PR); OR also possible.\\
Course example, high BP: 38/79 (48\%) smokers vs 86/401 (21\%) non-smokers. PR 2.24 (1.67 to 3.02); OR 3.39 (2.06 to 5.61).};
\node[m=8.6cm, below=4mm of cs] (cc) {\textbf{Case-control study}\\
Cases and controls are chosen by outcome, so the proportion with the outcome is set by the design. Risks cannot be estimated: \textbf{only the OR} is interpretable.};
\node[warn=8.6cm, below=4mm of cc] (rare) {\textbf{OR is close to RR only when the outcome is rare} (roughly under 10\%). With a common outcome the OR is further from 1: do not read an OR of 3.4 as ``3.4 times the risk''.};
\node[m=8.6cm, below=4mm of rare] (adj) {\textbf{Adjusting for a confounder}\\
Stratify and combine: \rc{mantelhaen.test(...)}, or logistic regression for an adjusted OR (aOR): see \dref{D4}.};
\end{tikzpicture}
\end{document}
